\chapter{Tổng kết}\label{chapter:conclusion}
\pagestyle{fancy}

\section{Kết quả đạt được}

Thực tập 1 đã hoàn thành ba mục tiêu đề ra ở Chương~\ref{chapter:introduction}.

\paragraph{Về tổng hợp bài báo.} Báo cáo đã trình bày một cách có hệ thống công trình FashionMV~\cite{yuan2026fashionmv}: bài toán Multi-View CIR và hiện tượng khiếm khuyết góc nhìn; pipeline ba giai đoạn xây dựng bộ dữ liệu FashionMV; mô hình ProCIR với ba cơ chế hội thoại hai lượt, căn chỉnh dựa trên chú thích và chuỗi suy luận với loại bỏ tiệm tiến, cùng giai đoạn SFT tuỳ chọn; và các kết quả thực nghiệm gồm so sánh với sáu mô hình hiện có và ablation 16 cấu hình. Công trình được đặt trong bối cảnh rộng hơn qua phần tổng quan lịch sử bài toán CIR (Chương~\ref{chapter:overview}) và cơ sở lý thuyết về Transformer, attention tuyến tính, học tương phản, CLIP và MLLM (Chương~\ref{chapter:background}). Ngoài ra, học viên đã đọc mã nguồn đánh giá để làm rõ những chi tiết giao thức không được nêu trong bài báo, và đưa ra các nhận xét phản biện về phương pháp (Mục~\ref{sec:method-critique}).

\paragraph{Về tái lập thực nghiệm.} Học viên đã xây dựng quy trình thay thế nguồn ảnh gốc không khả dụng bằng các bản sao trên HuggingFace, đạt độ phủ 100\% sản phẩm với DeepFashion và 64,4\% với Fashion200K, và chạy thành công mã đánh giá gốc với checkpoint công khai trên GPU miễn phí của Kaggle. Kết quả tái lập đạt R@5/R@10 = 74,42/85,25 trên toàn bộ 5.188 triplet DeepFashion và 87,93/94,53 trên mẫu 1.500 triplet Fashion200K. Các sai lệch so với bài báo -- thấp hơn trên DeepFashion, cao hơn trên Fashion200K -- đều được truy nguyên đến các nguyên nhân cụ thể và định lượng được: ảnh thay thế DeepFashion chỉ có độ phân giải $224 \times 224$ (cung cấp ít token thị giác hơn khoảng 5 lần so với ảnh gốc), và gallery Fashion200K bị thu nhỏ 4,5 lần do đánh giá trên mẫu con.

\paragraph{Về mô hình cơ sở.} Hai baseline zero-shot CLIP và FashionCLIP được đánh giá trên \emph{đúng} tập triplet và gallery của ProCIR. Thứ tự hiệu năng CLIP $<$ FashionCLIP $<$ ProCIR được quan sát nhất quán trên cả hai bộ dữ liệu, với khoảng cách lớn giữa ProCIR và baseline tốt nhất, khẳng định giá trị của kiến trúc CIR chuyên biệt trên nền MLLM so với hợp nhất muộn tuyến tính. Khảo sát trọng số hợp nhất và phân tích vai trò của sản phẩm nguồn trong gallery cung cấp thêm hiểu biết về giao thức đánh giá mà bài báo gốc chưa đề cập.

\paragraph{Về tái lập huấn luyện.} Sau khi có đủ ảnh của cả ba bộ dữ liệu, học viên cài đặt lại phần huấn luyện của ProCIR và huấn luyện hai cấu hình MT+Align và single-turn trên 25\% dữ liệu với LoRA (Mục~\ref{sec:train-repro}). Trên toàn bộ tập validation, MT+Align đạt trung bình R@5 = 72,47 so với 71,26 của single-turn, với cùng hình mẫu theo từng bộ dữ liệu như ablation của bài báo. Trên cùng bộ ảnh, checkpoint của tác giả đạt R@5 = 89,26 trên DeepFashion và 75,82 trên FashionGen, khớp với số công bố; kết quả trên DeepFashion kiểm chứng trực tiếp lập luận về độ phân giải nêu ở trên. Trên Fashion200K với gallery đầy đủ, checkpoint này đạt 67,71, thấp hơn bài báo khoảng 10 điểm; khoảng cách này giúp phát hiện một lỗi trong bước cắt ảnh Fashion200K của học viên, làm khoảng một phần tư số ảnh của bộ này bị cắt nhầm sang món đồ khác (Mục~\ref{sec:train-repro}).

\paragraph{Về kinh nghiệm thực tiễn.} Quá trình thực hiện cho thấy tái lập một MLLM có kiến trúc attention lai trên tài nguyên miễn phí là khả thi nhưng đòi hỏi nhiều lần điều chỉnh (cỡ mẫu, batch, thư viện kernel), và nhấn mạnh tầm quan trọng của việc ghi lại đầy đủ các chi tiết như nguồn ảnh, độ phân giải và cách dựng gallery khi báo cáo kết quả truy xuất.

\section{Những hạn chế tồn tại}

\begin{itemize}
    \item \textbf{Số liệu ở các Mục~\ref{sec:repro-results}--\ref{sec:discussion} đo trên bộ ảnh thay thế.} Các mục này dùng ảnh DeepFashion $224 \times 224$, một mẫu con của Fashion200K và không có FashionGen. Phần tái lập huấn luyện (Mục~\ref{sec:train-repro}) đã đánh giá cả ba bộ trên toàn bộ tập validation với ảnh đầy đủ; khi so sánh với bài báo nên dùng Bảng~\ref{tab:train-results}.
    \item \textbf{Số liệu Fashion200K ở Mục~\ref{sec:train-repro} đo trên ảnh cắt lỗi.} Khoảng một phần tư số ảnh Fashion200K, ở cả tập huấn luyện lẫn tập validation, bị cắt nhầm sang món đồ khác. Lỗi đã được sửa trong mã chuẩn bị dữ liệu, nhưng các số liệu Fashion200K chưa được đo lại và hai cấu hình chưa được huấn luyện lại trên ảnh đã sửa.
    \item \textbf{Fashion200K chỉ đánh giá trên mẫu con.} Mẫu 1.500/18.499 triplet với gallery 2.367 sản phẩm khiến kết quả trên bộ này không so sánh trực tiếp được với số liệu bài báo; ngoài ra tập triplet khả dụng (48,2\%) đã là một tập con chọn lọc theo độ phủ của nguồn ảnh thay thế.
    \item \textbf{Huấn luyện lại mới ở quy mô thu gọn.} Hai lần chạy dùng 25\% dữ liệu, LoRA, không có bước SFT và một seed; chúng được chấm ở hai bước huấn luyện khác nhau do tiêu chí chọn checkpoint trên tập dev còn thô. Trong 16 cấu hình ablation của bài báo, mới có hai cấu hình được kiểm chứng độc lập.
    \item \textbf{Baseline còn đơn giản.} Hai baseline là mô hình zero-shot với hợp nhất muộn; chưa đánh giá các baseline mạnh hơn trong bài báo như SPRC hay Qwen3-VL-Embedding do giới hạn tài nguyên.
    \item \textbf{Ảnh hưởng của việc thiếu kernel tối ưu chưa được đối chiếu trực tiếp.} Về nguyên tắc, chạy attention tuyến tính bằng triển khai tham chiếu chỉ ảnh hưởng tốc độ, không ảnh hưởng kết quả số học (ngoài sai số dấu phẩy động). Ở Mục~\ref{sec:train-repro}, kernel tối ưu chạy được và checkpoint của tác giả cho kết quả khớp bài báo trên hai bộ dữ liệu; tuy vậy hai triển khai chưa được so sánh trên cùng dữ liệu.
\end{itemize}

\section{Định hướng phát triển}

Từ các kết quả và hạn chế trên, học viên đề xuất các hướng phát triển sau cho giai đoạn Thực tập 2:

\begin{enumerate}
    \item \textbf{Hoàn thiện tái lập.} Đo lại Fashion200K trên ảnh đã cắt đúng và huấn luyện lại hai cấu hình trên dữ liệu đã sửa; chấm hai cấu hình ở cùng một bước; lặp lại với nhiều seed và, khi tài nguyên cho phép, trên toàn bộ dữ liệu huấn luyện.
    \item \textbf{Chuẩn hoá giao thức đánh giá.} Báo cáo song song kết quả có và không loại sản phẩm nguồn khỏi danh sách xếp hạng, và luôn báo cáo kèm kích thước gallery.
    \item \textbf{Phân tích cơ chế tổng hợp đa góc nhìn.} Khảo sát bản đồ attention của ProCIR để xem mô hình chú ý vào góc nhìn nào khi văn bản đề cập một chi tiết (ví dụ phần lưng), cũng như độ bền của mô hình khi thiếu góc nhìn, khi đảo thứ tự góc nhìn hoặc khi chỉ có một ảnh.
    \item \textbf{Đề xuất cải tiến.} Dựa trên các phân tích trên, nghiên cứu cơ chế tổng hợp đa góc nhìn có nhận biết góc nhìn (\emph{view-aware}) -- ví dụ gắn nhãn góc nhìn tường minh cho từng ảnh hoặc dùng trường \texttt{views\_involved} làm tín hiệu giám sát phụ -- nhằm giải quyết trực tiếp hơn hiện tượng khiếm khuyết góc nhìn, theo hướng đã nêu trong đề cương nghiên cứu. Vì dữ liệu công bố không kèm nhóm ứng viên bị loại trong quá trình xây dựng, kho mẫu âm khó mà đề cương dự kiến sẽ được khai thác từ embedding của mô hình đã huấn luyện ở Mục~\ref{sec:train-repro}.
\end{enumerate}

\section{Kế hoạch và tiến độ thực hiện}

Công việc được thực hiện từ ngày 07/09/2026 theo sáu nhóm: khảo sát tài liệu, nghiên cứu bài báo FashionMV, viết đề cương nghiên cứu, chuẩn bị dữ liệu và môi trường, thực nghiệm, và viết báo cáo (Hình~\ref{fig:timeline}). Phần tái lập huấn luyện ở Mục~\ref{sec:train-repro} được thực hiện tiếp sau các công việc trong hình, đến ngày 10/10/2026. Tiến độ được báo cáo hằng tuần với giảng viên hướng dẫn. Toàn bộ công việc do học viên thực hiện cá nhân dưới sự hướng dẫn của giảng viên.

\begin{figure}[h]
\centering
\includegraphics[width=\textwidth]{timeline_tt1.pdf}
\caption{Kế hoạch và tiến độ thực hiện Thực tập 1.}
\label{fig:timeline}
\end{figure}
